\section{Conclusion}
\label{sec:conclusion}

We introduced \ours{}, a minimal modification to standard deep
ensembling: train the same architecture with the same cosine schedule
but two different $T_\mathrm{max}$ endpoints, and average the resulting
softmax outputs. The two trajectories reach identical mean accuracy
yet occupy distinct loss basins whose decision boundaries disagree on
13\% of test samples---a free second source of ensemble diversity that
costs one extra training run and zero hyperparameter tuning.

Across 9 benchmarks, three modalities (EEG, images, text), and five
architectures (axial transformer, small CNN, ResNet-18, MLP,
DistilBERT), \ours{} delivers consistent gains over standard deep
ensembles at matched compute, including a new state of the art on
FACED 9-class EEG emotion classification (\facedacc{} BAcc, $+8.2\%$
over the EMOD AAAI 2026 baseline and $+1.1$~pp over a 10-seed
same-length deep ensemble) and a clean head-to-head sweep against SWA,
snapshot ensembles, and 10-seed deep ensembles across accuracy, NLL,
ECE, and Brier.

The mechanism is transparent and predictive. The low-LR tail of the
longer schedule drifts weights into a neighboring basin (confirmed by
mode-connectivity probes and weight-space PCA), producing
prediction-space diversity that weight averaging cannot reproduce. The
per-sample decomposition $q \cdot g$ quantitatively predicts the
ensemble gain ($0.131 \times 0.103 = 1.35$~pp, observed $1.50$~pp),
and the gain scales as a power law in the headroom $(1-p_{\mathrm{ind}})$
with $\alpha = 0.85$---so the trick pays off most on exactly the
capacity-bound tasks where ensembling is most needed.

A single integer change in the training schedule, a reliable gain
beyond a true 10-seed deep ensemble, and a predictive scaling law that
tells you in advance when to expect the gain: \ours{} is as close to a
``free lunch'' as deep ensembling currently offers. We expect it to
compose cleanly with any diversity axis that future work identifies.
